\section{Sex and Gender Non-Conformity in the Chinese Vinayas of the Different Schools}
\label{appendix1}

\subsection{Mahāsaṅghika Vinaya}

The term 非男非女 (neither male nor female) is only used by the \textit{paṇḍaka} to describe himself in this \textit{Vinaya}. This could be a literal translation of the term \textit{napuṁsaka} as in Vedic India this is an umbrella term, which includes the \textit{paṇḍaka}. The \textit{hijras} of India also refer to themselves with this term.

The term 二根 (i.e. 2 roots/faculties) is mentioned in passing as a question for \textit{Bhikkhu} ordination but without further explanation.\footnote{T22 1425 摩訶僧祇律 0413c02.} The term 黃門 (translated as `eunuch'\footnote{In the remainder of this appendix I will use the translation `eunuch' (in quotation marks) as the official translation of 黃門 according to the dictionary. I refute this translation as too narrow and probably erroneous in chapter \ref{pandakainchinese}.}) is also mentioned here without further explanation.

Another term used in \textit{Pārājika} 2\footnote{T22 1425 摩訶僧祇律 0244a24.} is 二形 (`two shapes'). There seems to be some confusion between two terms: 二根 (`two roots/faculties') and 二形 (`two shapes') that are sometimes used as synonyms in different places.

Other words we find in the \textit{Bhikkhunī} ordination procedure are those who have no breasts (無乳) or just one breast (一乳) or those who have no vagina (石女).\footnote{Dr. Hsiao-Lan Hu translates 石女 as `having no vagina' but the dictionary entry in \href{http://www.buddhism-dict.net/cgi-bin/xpr-ddb.pl?q=石女}{www.buddhism-dict.net} translates it as `barren' or `sterile'.} In this procedure it is mentioned that the candidate can proceed if she does not suffer from these conditions.\footnote{T22 1425 摩訶僧祇律 0472b05–0472b10.} We also find a question about whether or not she has a rectovaginal fistula (二道通).\footnote{The commentary T24 1462 善見律毘婆沙 0763 explains: ``二道者，穀道水道''--``The two tracts are the solid tract and the liquid tract.''}

\subsection{Dharmaguptaka Vinaya}
The regular list of persons not to be ordained is given, using the word 二形 (`two shapes'), translated by \cite{bodhi} as `hermaphrodite', while in other places in the \textit{Vinaya} it uses 二根 (`two roots/faculties'). 

After this list the \textit{Dharmaguptaka} \textit{Vinaya} adds the story of a monk and nun respectively who change sex/gender as mentioned in the \textit{Theravāda} \textit{Pārājika} 1. The Buddha concludes that they can simply go to the other Order and do not need to be expelled.\footnote{T22 1428 四分律 0813b15–0813b23. The commentary X55 0884: 表無表章栖翫記 0230c22–0231a09 explains that there is no need for re-ordination in this case.} Again, the word used here for sex/gender characteristics is 形 (i.e. form or shape). The next paragraphs list the case of a monk and nun respectively who changed sex/gender to become 男女二形 i.e. both male and female. The Buddha mentions that they have to be expelled.

The \textit{Dharmaguptaka} \textit{Vinaya} proceeds to list details of monks who have been castrated through various causes.\footnote{T22 1428 四分律 0813b25–0813c04.} Obviously these are not seen as falling under the same category as the above mentioned `eunuch'. When castration happens through an accident or even when it happens through \textit{karmic} causes, the monk in question can remain, but if he causes the castration intentionally himself he is expelled. Here the phrase is 截其 男根 (lit. cut off the male root/faculty).

While in the \textit{Mahāsaṅghika} \textit{Vinaya} one who is castrated (i.e. with the cutting off of the male root/faculty 男根) is seen as an impotent man and thus not fit for ordination, here this only matters when the action is voluntary and not accidental.

In the \textit{Bhikkhunī} ordination procedure we find the person with two faculties (二根) and the person with rectovaginal fistula (二道合). In the same sequence we find the word 道小, which is translated by \cite{bodhi} as `underdeveloped genitalia'.\footnote{T22 1428 四分律 0924c20.} Unlike in the \textit{Theravāda} \textit{Vinaya}, this condition would lead to disqualification from ordination.\footnote{The disqualification from ordination of the 二根 and 二道合 is confirmed in \textit{Bhikkhunī Vibhaṅga, Pācittiya} 166–167, T22 1428 四分律 0773c21–0774b03.} Further down a separate clause is added for those who have no breasts (無乳) or just one breast (一乳), who are equally barred from ordination.\footnote{T22 1428 四分律 0926c20–0926c21.}

\subsection{Mahīśāsaka Vinaya}
In the \textit{Bhikkhunī} ordination procedure we find another few terms in the questions asked during the ritual.\footnote{T22 1421 彌沙塞部和醯五分律 0187c21–0187c29.} It asks if a woman has no vagina (石女),\footnote{Translated elsewhere as `barren/sterile'.} if she is not a `eunuch' (黃門) and if her female genitals (root/faculty) are developed (女根具足). Here it is not specifically mentioned that somebody is barred from ordination if the answer is affirmative.

\subsection{Sarvāstivāda Vinaya}

In another part of the \textit{Vinaya} this term 二根 (two roots/faculties) is used next to the term 黃門 (`eunuch') but not in relation to ordination. \textit{Pārājika} 1 (just like the \textit{Pārājika} 1 of all the schools) mentions the existence of four kinds of offender: men, women, 黃門 (`eunuch') and 二根 (two roots/faculties). The same two words are used elsewhere in the \textit{Sarvāstivāda} \textit{Vinaya} while the word 不能男 (impotent) is only used in the list for those who cannot ordain.

The \textit{Bhikkhunī Khandhaka} goes more into detail about those who cannot ordain. The 二根 (two roots/faculties) is mentioned here.\footnote{T23 1435 0294a23–0294a28.} 

A similar list of questions is asked of female candidates for ordination as with the other schools. The candidate is to be asked if she has underdeveloped genitalia (女根小) (lit. small female root/faculty), has no breasts (無乳) or just one breast (一乳) and if she is sterile (是不能產). It seems however that regardless of the answer, the candidate is not barred from ordination.\footnote{T23 1435 0332b11–0332b22.}

With regards to the five types of 黃門, the Chinese commentarial texts merely add that these people cannot ordain because they have difficulty keeping the precepts.\footnote{T85 2792 毘尼心 0667b25–0667b26.}
